\documentclass[11pt,a4paper]{article}
\usepackage[margin=25mm]{geometry}
\usepackage{amsmath,amssymb,amsthm,hyperref}
\setlength{\emergencystretch}{2em}
\newtheorem{lemma}{Lemma}
\newtheorem{theorem}[lemma]{Theorem}
\newtheorem{corollary}[lemma]{Corollary}
\newcommand{\gen}[1]{\langle#1\rangle}
\title{The orbit-incidence deficit between separated bands}
\author{Human Compression project: L044 continuation}
\date{}
\begin{document}\maketitle
\begin{lemma}[Two generating subgroups]\label{incidence}
Let $G=\gen{P,R}$ act transitively on a finite set $\Omega$
of size $n>1$. Suppose $P,R$ have nonempty supports
$Y=\operatorname{supp}(P)$ and $X=\operatorname{supp}(R)$,
and respectively $p,q$ nontrivial orbits. Let $H\leq P\cap R$,
put $W=X\cap Y$, and let $o_W(H)$ count all $H$-orbits in
$W$, including singleton orbits. Then
\[
 X\cup Y=\Omega,\qquad o_W(H)\geq p+q-1.                 \tag{1}
\]
In particular every nontrivial $P$- or $R$-orbit meets $W$.
\end{lemma}
\begin{proof}
Form the bipartite multigraph whose vertices are all $P$-orbits
and all $R$-orbits, with one edge for each point of $\Omega$
joining the two orbits that contain it. It is connected:
a path through successive orbit nodes is exactly a chain of
moves by $P$ and $R$, and these generate the transitive group $G$.
A point fixed by both subgroups would be fixed by $G$, so no
such point exists and $X\cup Y=\Omega$.

A singleton orbit node is a leaf. Delete all such nodes and their
incident edges. The remaining nodes are the $p+q$ nontrivial
orbits, and the remaining edges are precisely the points of $W$.
This graph is connected: a path between two retained nodes cannot
have a leaf as an internal node. Both sorts of retained nodes
exist by the nonempty-support assumption. Thus no retained node
is isolated. Since $H\leq P\cap R$, each of its orbits in $W$
is contained in one intersection of a $P$-orbit and an $R$-orbit.
Replace the parallel edges belonging to each such $H$-orbit by
one edge. This preserves connectivity, and gives a connected
multigraph with $p+q$ vertices and $o_W(H)$ edges. The elementary
spanning-tree bound proves (1).
\end{proof}

\begin{theorem}[Separated interval with outside budgets]\label{deficit}
Let $G=\gen{\rho_0,\ldots,\rho_{r-1}}$ be a transitive string
involution group whose maximal omissions are intransitive.
Let $0\leq h<g\leq r-1$, and put
\[
 P=G_{<g},\quad R=G_{>h},\quad H=G_{\{h+1,\ldots,g-1\}},
 \quad Y=\operatorname{supp}(P),\quad X=\operatorname{supp}(R).
\]
The empty interval means $H=1$. Suppose each colour strictly
between $h$ and $g$ has two ambient-eligible edges, and
\[
 2h\leq n-|X|-1,\qquad 2(r-1-g)\leq n-|Y|-1.             \tag{2}
\]
If $p,q$ are the numbers of nontrivial $P,R$-orbits, then
\[
                         2r\leq n+3-p-q.                \tag{3}
\]
Consequently, if both groups are intransitive on their supports,
then $r\leq\lfloor(n-1)/2\rfloor$.
\end{theorem}
\begin{proof}
The index sets of $P,R$ cover every generator because $g>h$,
so $G=\gen{P,R}$. Both supports are nonempty. Moreover
$H\leq P\cap R$; no intersection property is needed for this
inclusion. Set $W=X\cap Y$ and $V=\operatorname{supp}(H)$.
Then $V\subseteq W$.

For a nonempty middle interval, $H$ lies in a maximal omission
and is therefore intransitive on $\Omega$. The complete accepted
ambient interval theorem gives, with $c$ its number of nontrivial
orbits,
\[
 2(g-h-1)\leq |V|-c=|W|-o_W(H)\leq |W|-p-q+1.           \tag{4}
\]
The equality includes every $H$-fixed point in $W$; such points
are not silently discarded. The last inequality is Lemma~\ref{incidence}.
If $g=h+1$, then $H=1$, $V=\varnothing$ and
$|W|-o_W(H)=0$, so (4) remains true with left side zero.
Adding (2) and (4), and using $|X|+|Y|=n+|W|$, gives
\[
\begin{split}
2r&=2h+2(g-h-1)+2(r-1-g)+4\\
&\leq (n-|X|-1)+(|W|-p-q+1)+(n-|Y|-1)+4\\
&=n+3-p-q.
\end{split}
\]
When both groups are intransitive on their supports, $p,q\geq2$,
so $2r\leq n-1$ as required.
\end{proof}

\begin{corollary}[Oriented two-bridge conclusion]\label{bridge}
Use the standing $G\leq A_n$, $n\geq12$, full string-IP hypotheses
of the accepted nested-bridge theorem, including its nontrivial
outside regions and prefix/suffix orientation. Suppose $h$ is
last shared across the first cut and $g$ first shared across the
second, with $i<h<j$, $g<j$, and the stated pure outside actions.
Suppose also that every colour strictly between $h$ and $g$
when $g>h$ has two ambient-eligible edges, and the accepted
outside-support bounds apply in both directions. Under the same
smaller-degree alternating induction as the accepted transitive-tail
alternative, either
\[
 r\leq\lfloor(n-1)/2\rfloor,
\]
or an actual first or last three-generator subgroup is one of the
two Table~2 six-point $A_5$ markings, in the appropriate orientation.
\end{corollary}
\begin{proof}
Assume the rank bound fails. The accepted nested-bridge count gives
$g\geq h$, and its accepted interior and endpoint touching
exclusions give $g>h$. The assumed support bounds are exactly (2),
with the reversed prefix index $r-1-g$.
If both $P,R$ are intransitive on their supports,
Theorem~\ref{deficit} contradicts rank failure. Otherwise apply
the complete accepted transitive-tail alternative to the first
cut if $R$ is transitive, or in the literal reversed orientation
to the second cut if $P$ is transitive. Its exact unique-crossing,
shared-colour and orientation hypotheses are included here.
It supplies precisely the stated actual end subgroup.
\end{proof}

\paragraph{Attribution and remaining source routing.}
This replaces the last orbit count in Cameron--Fernandes--Leemans--Mixer,
\emph{Highest rank of a polytope for $A_n$}, arXiv:1605.09173v1,
Proposition~5.28, physical page~37. The source uses the support overlap
and intransitivity strategy. The explicit orbit-incidence multigraph
and the $p+q-1$ deficit are the present replacement argument.
They retain singleton intersection orbits and the corrected reversed
index, and avoid assuming that two intransitive tails yield only a
two-orbit middle deficit.

All ambient-interval, nested-bridge, touching and transitive-tail
proofs and their complete required programs remain paid. No new
finite computation or formal declaration is needed. The corollary
is deliberately conditioned on the actual cut orientations and
support bounds: their complete routing from the source's maximal
choice, and its more specific choice of end, remain separate.
This packet does not claim whole L044 recovery, new original,
question or priority credit.
\end{document}
